\chapter{Results and Discussion}
\label{ch:results-discussion}

%% ── 5.1 ─────────────────────────────────────────────────────────────────────
\section{Evaluation Methodology}
\label{sec:eval-methodology}

\subsection{Metrics}
% R² (coefficient of determination): primary metric, directly comparable to prior project
% RMSE (root mean square error) in dB: interpretable error magnitude
% MAE (mean absolute error) in dB: robust to outliers
% Communication cost: bytes per device per day

\subsection{Experimental Protocol}
% Same architecture for centralized and federated (Dense 5→8→1)
% Same preprocessing, same normalization
% Same held-out test set (20% of data, stratified by device)
% Centralized: trained on ALL training data
% Federated: trained on device-partitioned data via FedAvg simulation
% This ensures the comparison is fair — only the training paradigm differs

%% ── 5.2 ─────────────────────────────────────────────────────────────────────
\section{Centralized Baseline Results}
\label{sec:centralized-results}

% R² = [XX], RMSE = [XX] dB, MAE = [XX] dB
% COMPARE WITH PRIOR PROJECT:
%   "The centralized neural network achieves R² = [XX], compared to the prior
%    project's XGBoost R² ≈ 0.93. The difference of [XX] is attributable to
%    the simpler architecture (57 params vs thousands in XGBoost), which is
%    a necessary trade-off for microcontroller deployment."

% FIGURES:
% - Training/validation loss curves over 50 epochs
% - Residual plot: predicted vs actual path loss (scatter)
% - Error histogram

%% ── 5.3 ─────────────────────────────────────────────────────────────────────
\section{Federated Learning Results}
\label{sec:fl-results}

\subsection{Convergence Analysis}
% FIGURE: R² vs communication round number (main result plot)
% How many rounds to reach 90% of centralized R²?
% How many rounds to converge?
% Compare with Torres Sanchez: they found 20 epochs × 4 rounds optimal

\subsection{Impact of Non-IID Data}
% FIGURE: IID random split vs non-IID device-based split — convergence curves
% Which devices' data deviates most from the global distribution?
% Per-client local model R² vs global model R²

\subsection{Hyperparameter Sensitivity}
% TABLE/FIGURES showing impact of:
%   - Local epochs: 1, 3, 5
%   - Minimum clients per round: 2, 3, 4, 6
%   - Total communication rounds: 5, 10, 20, 50
% Identify optimal configuration
% Compare with Torres Sanchez finding: "20 epochs × 4 rounds"

\subsection{Final Federated Model}
% Best configuration: [XX] rounds, [XX] local epochs, [XX] min clients
% R² = [XX], RMSE = [XX] dB, MAE = [XX] dB

%% ── 5.4 ─────────────────────────────────────────────────────────────────────
\section{Three-Way Performance Comparison}
\label{sec:three-way-comparison}

% THE KEY TABLE OF THE THESIS:
\begin{table}[H]
  \centering
  \caption{Three-way comparison: prior centralized XGBoost, centralized NN, and federated NN.}
  \label{tab:three-way}
  \begin{tabular}{lccc}
    \toprule
    \textbf{Metric} &
    \textbf{XGBoost} &
    \textbf{Centralized NN} &
    \textbf{Federated NN} \\
    & \textbf{(Prior Project)} & \textbf{(This Thesis)} & \textbf{(This Thesis)} \\
    \midrule
    $R^2$                          & $\approx 0.93$ & [XX]   & [XX]   \\
    RMSE (dB)                      & [XX]           & [XX]   & [XX]   \\
    MAE (dB)                       & [XX]           & [XX]   & [XX]   \\
    Model parameters               & $>$1000        & 57     & 57     \\
    Model size                     & $>$100\,KB     & 2.7\,KB & 2.7\,KB \\
    Runs on MCU?                   & No             & Yes    & Yes    \\
    Raw data leaves device?        & Yes            & Yes    & No     \\
    Adapts post-deployment?        & No             & No     & Yes    \\
    Daily bandwidth / node         & $\approx$25.9\,KB & $\approx$25.9\,KB & $\approx$2.4\,KB \\
    \bottomrule
  \end{tabular}
\end{table}

%% ── 5.5 ─────────────────────────────────────────────────────────────────────
\section{Communication Efficiency}
\label{sec:comm-efficiency}

% DETAILED CALCULATION:
\begin{table}[H]
  \centering
  \caption{Daily communication cost comparison.}
  \label{tab:bandwidth}
  \begin{tabular}{lrrr}
    \toprule
    & \textbf{Old System} & \textbf{New System} & \textbf{Reduction} \\
    \midrule
    Payload per message      & 18\,B       & 8\,B (status) / 52\,B (FL) & --- \\
    Messages per day / node  & 1440        & 288 (status) + 1 (FL)      & --- \\
    Bytes per day / node     & 25,920\,B   & $\approx$2,404\,B          & $\approx$11$\times$ \\
    Bytes per day / 6 nodes  & 155,520\,B  & $\approx$14,424\,B         & $\approx$11$\times$ \\
    \bottomrule
  \end{tabular}
\end{table}

% COMPARE WITH TORRES SANCHEZ:
% Their 1.39 KB model required 7–28 messages per round per client
% Our 52-byte update requires exactly 1 message per round per client
% Our model is ~27× smaller, so communication cost is dramatically lower

% EU868 DUTY CYCLE ANALYSIS:
% At DR0 (SF12, 51B payload): airtime per message ≈ 2.5 s
% 1% duty cycle: max ~36 messages/hour
% Old system (1440/day = 60/hour): EXCEEDS duty cycle at DR0
% New system (289/day = 12/hour): comfortably within duty cycle

%% ── 5.6 ─────────────────────────────────────────────────────────────────────
\section{Hardware Feasibility}
\label{sec:hardware-feasibility}

% COMPILE AND REPORT:
\begin{table}[H]
  \centering
  \caption{MKR WAN 1310 resource usage.}
  \label{tab:hw-resources}
  \begin{tabular}{lrrl}
    \toprule
    \textbf{Resource} & \textbf{Used} & \textbf{Available} & \textbf{Utilization} \\
    \midrule
    Flash              & [XX]\,KB  & 256\,KB   & [XX]\% \\
    SRAM               & [XX]\,KB  & 32\,KB    & [XX]\% \\
    Tensor arena       & [XX]\,B   & 8,192\,B  & [XX]\% \\
    Inference latency  & [XX]\,ms  & ---       & --- \\
    \bottomrule
  \end{tabular}
\end{table}

%% ── 5.7 ─────────────────────────────────────────────────────────────────────
\section{Answering the Research Questions}
\label{sec:answer-rqs}

% Answer EACH RQ explicitly with specific numbers:

% RQ1: Can federated TinyML achieve comparable R² to centralized?
%   "The federated model achieves R² = [XX] compared to centralized R² = [XX],
%    a relative difference of [X]%. This demonstrates that federated learning
%    can maintain [acceptable/strong] prediction accuracy while distributing
%    the training process across edge devices."

% RQ2: How does non-IID affect convergence?
%   "Non-IID partitioning by device location [slowed/did not significantly slow]
%    convergence compared to IID random splitting. The global model required
%    [XX] rounds under non-IID vs [XX] rounds under IID to reach 90% of
%    centralized R². This aligns with Torres Sanchez et al.'s finding that
%    FL is feasible even with non-IID and unbalanced client data."

% RQ3: What bandwidth reduction?
%   "The federated approach reduces daily per-node communication from
%    25,920 bytes to approximately 2,404 bytes, an 11× reduction."

% RQ4: Does it fit on MKR WAN 1310?
%   "The full pipeline occupies [XX]% of Flash and [XX]% of SRAM,
%    confirming feasibility on the Cortex-M0+ platform."

%% ── 5.8 ─────────────────────────────────────────────────────────────────────
\section{The Centralized-to-Federated Trade-Off}
\label{sec:tradeoff}

% What do we lose? R² drops by [X]%
% What do we gain? 11× bandwidth reduction, privacy, adaptability
% Is the trade-off worthwhile?
% "For long-running LoRaWAN deployments under duty-cycle constraints, the
%  marginal accuracy loss is acceptable given the substantial communication
%  and privacy benefits."

%% ── 5.9 ─────────────────────────────────────────────────────────────────────
\section{Environmental Influence on Path Loss}
\label{sec:env-influence}

% Feature importance analysis:
%   Which features contribute most to path loss prediction?
%   Physical interpretation: humidity → RF absorption, CO2 → occupancy proxy
% Does the relationship differ by location?
%   Per-client model analysis: some rooms may have stronger env-PL coupling

%% ── 5.10 ─────────────────────────────────────────────────────────────────────
\section{Comparison with Torres Sanchez et al.}
\label{sec:comparison-torres}

% DIRECT SIDE-BY-SIDE (reference Table 3.1):
% They: anomaly detection, Flower, 4 virtual clients, 77K messages, F1=94.77%
% We: path loss regression, custom FedAvg, 6 virtual clients, 1.7M rows, R²=[XX]

% SHARED FINDINGS:
% 1. FL achieves comparable performance to centralized under LoRaWAN constraints
% 2. Model size must be minimized — they needed 7+ messages, we need exactly 1
% 3. Round/epoch configuration matters

% OUR ADVANTAGES:
% 1. Real measurement data from physically deployed nodes (vs machinery sensors)
% 2. Non-IID split by actual physical location (vs per-machine-type)
% 3. Smaller model: 52 bytes vs 1.39 KB → fits in single LoRaWAN message
% 4. Hardware prototype on real MKR WAN 1310 with TFLite Micro

%% ── 5.11 ─────────────────────────────────────────────────────────────────────
\section{Comparison with the Prior Project}
\label{sec:comparison-prior}

% Three-way comparison narrative (reference Table 6.3):
% Prior XGBoost: R²≈0.93, but huge model, centralized, no privacy, no adaptation
% Centralized NN: R²=[XX], small model, still centralized
% Federated NN: R²=[XX], small model, distributed, private, adaptive
% "The progression from XGBoost to federated NN represents a deliberate
%  trade-off: accepting a modest reduction in R² in exchange for deploying
%  intelligence directly at the edge."

%% ── 5.12 ─────────────────────────────────────────────────────────────────────
\section{Limitations}
\label{sec:limitations}

\begin{enumerate}
  \item \textbf{Simplified local training.} The on-device training uses a separate
    weight proxy array rather than true backpropagation through the TFLite model.
    This limits the fidelity of on-device learning compared to full gradient descent.

  \item \textbf{Simulation-based FL evaluation.} While the data is real, the federated
    process is simulated on a~PC. A multi-week live deployment with all six nodes
    communicating through actual LoRaWAN would provide stronger validation.
    Torres Sanchez et al.\ also relied on simulation, confirming this as standard
    practice, but live validation remains desirable.

  \item \textbf{Single building, single gateway.} Results may not generalize to
    multi-building or outdoor deployments with different propagation characteristics.

  \item \textbf{Proxy labels for on-device training.} During deployment, devices use
    \gls{pdr} as a proxy for link quality rather than direct path loss measurement.
    The accuracy of this proxy has not been formally validated.

  \item \textbf{Fixed model architecture.} Only one architecture (Dense 5$\to$8$\to$1)
    was evaluated. Deeper or wider models might improve accuracy but would increase
    communication cost.
\end{enumerate}
